\chapter{Склад та функціональна структура ERP/1}

\section{Концепція побудови екосистеми ІКС}
Інформаційно-комунікаційні системи (ІКС) платформи «ERP/1» розроблені як комплекс взаємоінтегрованих прикладних рішень для автоматизації державного та комерційного секторів. Усі продукти функціонують на базі єдиного технологічного стеку Erlang/OTP, що гарантує високу масштабованість, стійкість до відмов та безпеку персональних даних. Нижче наведено детальний опис та вимоги до кожного з 10 ключових продуктів ІКС.

---

\section{1. Освіта (LMS / EMS)}
\subsection{Опис та призначення}
Автоматизована інформаційна система управління освітнім процесом, навчальною діяльністю та адмініструванням закладів освіти (згідно з ISO 21001) з підтримкою очної, дистанційної та змішаної форм навчання.
\subsection{Ключові функціональні модулі}
\begin{itemize}
    \item \textbf{Навчальний процес}: електронні курси (SCORM), розклад занять, інтерактивні журнали, та фіксація відвідуваності.
    \item \textbf{Дистанційне навчання}: вбудовані відеоконференції, інтерактивні тести для контролю знань.
    \item \textbf{Контингент}: особові справи учнів та студентів, рух контингенту та академічна мобільність.
    \item \textbf{Кадри та атестація}: особові справи викладачів, облік підвищення кваліфікації, та розрахунок KPI.
    \item \textbf{Інтеграція}: автоматична взаємодія з ЄДЕБО через захищений протокол MQTT та PKI-інфраструктуру.
\end{itemize}

---

\section{2. Здоров'я (MIS / HL7)}
\subsection{Опис та призначення}
Медична інформаційна система (МІС) для автоматизації медичних послуг та управління медичною інформацією в електронному вигляді згідно з міжнародним стандартом HL7 FHIR.
\subsection{Ключові функціональні модулі}
\begin{itemize}
    \item \textbf{Електронна медична картка (EHR)}: історія хвороб, рецепти, направлення та протоколи лікування.
    \item \textbf{Реєстр пацієнтів}: ведення бази пацієнтів та ідентифікація через демографічні дані.
    \item \textbf{Управління прийомами}: розклад лікарів, електронна черга та онлайн-запис пацієнтів.
    \item \textbf{Взаємодія з eHealth}: двостороння інтеграція з центральною базою даних eHealth України для верифікації декларацій та звітів для НСЗУ.
\end{itemize}

---

\section{3. Документи (СЕД / CRM)}
\subsection{Опис та призначення}
Автоматизація роботи з електронними документами та впровадження електронного документообігу (СЕД) згідно з Державною інструкцією з діловодства №40 від 2024 року.
\subsection{Ключові функціональні модулі}
\begin{itemize}
    \item \textbf{Реєстрація кореспонденції}: вхідні, вихідні, внутрішні документи та звернення громадян.
    \item \textbf{Управління життєвим циклом}: BPE-процеси погодження, підписання (КЕП) та накладання резолюцій.
    \item \textbf{Сховище документів}: організація ієрархічної структури папок на базі KVS/Mnesia з контролем доступу ABAC.
\end{itemize}

---

\section{4. Облік (ERP / Accounting)}
\subsection{Опис та призначення}
Комплексна автоматизація бухгалтерського обліку, кадрового діловодства та розрахунку заробітної плати для державних установ та приватних підприємств відповідно до законодавства України.
\subsection{Ключові функціональні модулі}
\begin{itemize}
    \item \textbf{Головна книга}: план рахунків, подвійний запис, облік ТМЦ, ОЗ, ПДВ, розрахунки з контрагентами.
    \item \textbf{Кадровий облік}: штатні розписи, накази по персоналу, особові картки П-2, та табель обліку робочого часу.
    \item \textbf{Заробітна плата}: нарахування окладів, премій, лікарняних, розрахунок ЄСВ та ПДФО, генерація податкової звітності.
\end{itemize}

---

\section{5. Склад (WMS)}
\subsection{Опис та призначення}
Автоматизована система управління складським господарством, матеріально-технічним забезпеченням та логістичними ланцюгами постачання.
\subsection{Ключові функціональні модулі}
\begin{itemize}
    \item \textbf{Адресне зберігання}: динамічний розподіл складських зон, стелажів та осередків.
    \item \textbf{Операційний облік}: приймання, переміщення, комплектація, відвантаження ТМЦ та інвентаризація.
    \item \textbf{Штрихкодування}: інтеграція з терміналами збору даних (ТЗД) та підтримка кодів маркування GS1.
\end{itemize}

---

\section{6. Реєстри (Registry Platform)}
\subsection{Опис та призначення}
Універсальна low-code платформа для створення, ведення та інтеграції державних, відомчих чи корпоративних інформаційних реєстрів будь-якого масштабу.
\subsection{Ключові функціональні модулі}
\begin{itemize}
    \item \textbf{Конструктор схем}: візуальне проектування схем даних та автоматична генерація таблиць у сховищі.
    \item \textbf{Трасування змін}: фіксація кожного запису у вигляді незмінного ланцюжка транзакцій з накладанням КЕП.
    \item \textbf{Шина Трембіта}: вбудований адаптер для взаємодії з національним простором електронної взаємодії.
\end{itemize}

---

\section{7. Істотність (AI / NLP)}
\subsection{Опис та призначення}
Локальна ШІ-інфраструктура для інтелектуального аналізу, семантичного пошуку та автоматичної оцінки істотності інформації у документах та реєстрах без вивантаження даних у хмару.
\subsection{Ключові функціональні модулі}
\begin{itemize}
    \item \textbf{Семантичний аналіз}: класифікація документів, пошук за змістом та виявлення суперечностей.
    \item \textbf{CUDA-інференс}: робота з моделями Llama-3 (8B) та Phi-3 (14B) на споживчих GPU (16GB VRAM) за технологіями PagedAttention та Continuous Batching.
\end{itemize}

---

\section{8. Аналітика (DWH / OLAP)}
\subsection{Опис та призначення}
Високопродуктивна аналітична інфраструктура DWH для векторизованої обробки великих обсягів даних on-premises на базі вбудованого рушія DuckDB/MonetDB.
\subsection{Ключові функціональні модулі}
\begin{itemize}
    \item \textbf{Обробка Parquet}: генерація стовпчикових звітів безпосередньо з локального S3 сховища.
    \item \textbf{C99 NIF інтеграція}: швидке виконання OLAP запитів через Dirty CPU Schedulers планувальника Erlang.
\end{itemize}

---

\section{9. Проєкти (PM)}
\subsection{Опис та призначення}
Локальна заміна Atlassian Jira та Confluence на базі Erlang/OTP, розроблена для управління життєвим циклом проектів розробки ПЗ та ведення баз знань.
\subsection{Ключові функціональні модулі}
\begin{itemize}
    \item \textbf{Agile-управління}: Scrum/Kanban дошки, спринти, беклоги та BPE BPMN workflow для задач.
    \item \textbf{LaTeX Wiki}: ведення всієї документації проекту у форматі LaTeX з автоматичною генерацією ТВ, ТЗ та ТРП видавничої якості у PDF.
\end{itemize}

---

\section{10. Інциденти (ITSM)}
\subsection{Опис та призначення}
Система управління ІТ-послугами (ITSM) відповідно до рекомендацій бібліотеки найкращих практик ITIL для реєстрації та вирішення інцидентів, запитів та проблем.
\subsection{Ключові функціональні модулі}
\begin{itemize}
    \item \textbf{Service Desk}: прийом та класифікація звернень користувачів.
    \item \textbf{SLA-контроль}: автоматичне відстеження часових лімітів вирішення інцидентів.
    \item \textbf{База знань відомих помилок (KEDB)}: інтеграція з Вікі для швидкого вирішення типових збоїв.
\end{itemize}
